\documentclass[main]{subfiles}


\title[Dynamic Approaches]{Dynamic Approaches}
\subtitle{Background: Reinforcement Learning}
\author[André Biedenkapp]{André Biedenkapp}
\date{}

\titlebackground*{images/AutoML_LearningToLearn2}

%\institute{}
%\date{}

\begin{document}
	
	\maketitle
%----------------------------------------------------------------------
%-----------------------------------------------------------------------
%----------------------------------------------------------------------
\begin{frame}{Components of RL Problems}
    \centering
    \textbf{Goal:} An agent has to learn how to behave to maximize a reward
    \begin{tikzpicture}[node distance=4cm]
        %PreProcessing
        \node (Agent) [activity,
        minimum height=.9cm] {Agent};
        \onslide<2->{
        \node (Algo) [activity, right of=Agent, xshift=3cm,
        minimum height=.9cm] {Environment};}
        
        \begin{pgfonlayer}{background}
        \path (Agent -| Agent.west)+(-0.12,1.75) node (resUL) {};
        \path (Algo.east |- Algo.south)+(0.125,-1.75) node (resBR) {};
        % Top level
        \path [rounded corners, draw=black!50, fill=white] (resUL) rectangle (resBR);
        
        \end{pgfonlayer}
        \onslide<3->{
        \draw[myarrow, dashed] ($(Algo.south)+(-0.25, 0)$) -- ($(Algo.south)+(-0.25, -0.5)$) -- ($(Agent.south)+(0.25, -0.5)$) node [above,pos=0.5] {state $s_{t}$} -- ($(Agent.south)+(0.25, 0)$);}
        \onslide<4->{
        \draw[myarrow] (Agent.north) -- ($(Agent.north)+(0.0,+0.5)$) -- ($(Algo.north)+(0.0,+0.5)$) node [above,pos=0.5] {action $a_{t+1}$} node [below,pos=0.5] {} -- (Algo.north);}
        \onslide<5->{
        \draw[myarrow] ($(Algo.south)+(0.25, 0)$) -- ($(Algo.south)+(0.25, -0.95)$) -- ($(Agent.south)+(-0.25, -0.95)$) node [below,pos=0.5] {reward $r_{t}$} -- ($(Agent.south)+(-0.25, 0)$);}
    \end{tikzpicture}\\
    \phantom{Policy $\policy$: mapping from observations to actions}\\
    \phantom{Agent: entire decision making entity that contains the policy}
\end{frame}
%-----------------------------------------------------------------------
%----------------------------------------------------------------------
\begin{frame}{Agent vs. Policy}
    \centering
    \textbf{Goal:} An agent has to learn how to behave to maximize a reward
    \begin{tikzpicture}[node distance=4cm]
        %PreProcessing
        \node (Agent) [activity,
        minimum height=.9cm] {Policy $\pi$};
        
        \node (Algo) [activity, right of=Agent, xshift=3cm,
        minimum height=.9cm] {Environment};
        
        \begin{pgfonlayer}{background}
        \path (Agent -| Agent.west)+(-0.12,1.75) node (resUL) {};
        \path (Algo.east |- Algo.south)+(0.125,-1.75) node (resBR) {};
        
        % Top level
        \path [rounded corners, draw=black!50, fill=white] (resUL) rectangle (resBR);
        \end{pgfonlayer}
        
        \draw[myarrow] (Agent.north) -- ($(Agent.north)+(0.0,+0.5)$) -- ($(Algo.north)+(0.0,+0.5)$) node [above,pos=0.5] {action $\pi(s_t)\to a_{t+1}$} node [below,pos=0.5] {} -- (Algo.north);
        \draw[myarrow, dashed] ($(Algo.south)+(-0.25, 0)$) -- ($(Algo.south)+(-0.25, -0.5)$) -- ($(Agent.south)+(0.25, -0.5)$) node [above,pos=0.5] {state $s_{t}$} -- ($(Agent.south)+(0.25, 0)$);
        \draw[myarrow] ($(Algo.south)+(0.25, 0)$) -- ($(Algo.south)+(0.25, -0.95)$) -- ($(Agent.south)+(-0.25, -0.95)$) node [below,pos=0.5] {reward $r_{t}$} -- ($(Agent.south)+(-0.25, 0)$);
        
    \end{tikzpicture}\\\pause
    Policy $\policy$: mapping from observations to actions\\
    Agent: entire decision making entity that contains the policy
\end{frame}
%-----------------------------------------------------------------------
%----------------------------------------------------------------------
\begin{frame}{What Even is a State?}
    \begin{itemize}
        \item We constantly observe our surroundings and our internal feelings
        \item Senses provide different forms of observations:
        \begin{itemize}
            \item images
            \item sound
            \item feeling of touch
            \item feeling of acceleration
            \item feeling of balance
            \item \ldots
        \end{itemize}
    \end{itemize}
    It is not always possible to observe the full state.\\
    For simplicities sake, in the rest of the lecture we assume that we can observe all information that is relevant for an agent to solve a task.\footnote{See chapter 8 of \liturl{Puterman and Chan, 2026}{https://github.com/martyput/MDP_book/blob/master/08\%20Chapter/Book___chapter08.pdf} for a detailed discussion of the partial observation setting.}
\end{frame}
%-----------------------------------------------------------------------
%----------------------------------------------------------------------
\begin{frame}{Actions}
    \begin{itemize}
        \item An action will induce some change the state of the environment\pause
        \item Types of actions:
        \begin{description}
            \item[continuous] The value domain is continuous (and often bounded by some range, e.g., $\left[0,1\right]$)\pause\\
            Examples ...
            \begin{itemize}
                \item ... from robotics: velocities, angles, probabilities
                \item ... video-game playing: controller inputs
                \item ... from AutoML: learning rates, mutation rates
            \end{itemize}
            \item[discrete] Choose from a set of possible options\pause\\
            Examples ...\\
            \begin{itemize}
                \item ... from robotics: run sub-routine, interrupts
                \item ... video-game playing: keyboard inputs
                \item ... from AutoML: mutation operator selection, heuristic selection
            \end{itemize}
        \end{description}
    \end{itemize}
\end{frame}
%-----------------------------------------------------------------------
%----------------------------------------------------------------------
\begin{frame}{Transitions}
    \begin{itemize}
        \item Given a state $s$ and action $a$, in which state do we end up?\pause
        \item Can be either deterministic or non-deterministic
        \begin{itemize}
            \item Board games like Go or Chess are deterministic
            \item Games with randomized events (e.g., many card games) or robotics (control over the robot is not perfect) are non-deterministic
        \end{itemize}\pause
        \item Challenges
        \begin{itemize}
            \item Was the action responsible for the stochasticity (epistemic uncertainty) or is it inherent in the environment (aleatoric uncertainty)?
            \item Stochastic environments are more challenging since you need to account for a different notion of reproducibility!
        \end{itemize}
        \item Note: The transition function does not need to be known in order to learn well performing policies!
    \end{itemize}
\end{frame}
%-----------------------------------------------------------------------
%----------------------------------------------------------------------
\begin{frame}{Rewards}
    \begin{itemize}
        \item Rewards provide crucial feedback on wether the outcome of an action was ``good'' or ``bad''\pause
        \item Either immediate (or dense) reward: We directly get a reward signal after each transition\pause
        \item Or delayed (or sparse) reward: We have to wait for some states to observe the outcome of a previous choice or whole sequence of choices
    \end{itemize}\pause
    We need to balance reasoning about immediate gratification vs long-term ramifications!\\
    Particularly difficult: Temporal credit assignment during learning (see, e.g., \liturl{Rajan et al. 2023}{https://www.jair.org/index.php/jair/article/view/14314/26946})
\end{frame}
%-----------------------------------------------------------------------
%----------------------------------------------------------------------
\begin{frame}{Recall: Learning is an Iterative Process} 

    \begin{columns}
    \column{0.5\textwidth}
    Machine learning is \emph{iterative}
    %\pause
    \begin{enumerate}
        \item Initialize model
    %\pause
        \item Fit model to data
    %\pause
        \item Observe performance / compute loss
    %\pause
        \item Adapt model based on loss
    %\pause
        \item Repeat 2 - 4 until learning budget is exhausted
    \end{enumerate}\pause
    \column{0.5\textwidth}
    Reinforcement learning is \emph{iterative}!
    \pause
    \begin{enumerate}
        \item Initialize agent
    \pause
        \item \alert{Generate observations}
    \pause
        \item Observe performance / compute loss
    \pause
        \item Adapt policy based on loss
    \pause
        \item Repeat 2 - 4 until learning budget is exhausted
    \end{enumerate}
    \end{columns}
\end{frame}
%-----------------------------------------------------------------------
%----------------------------------------------------------------------
\begin{frame}{Q-Learning}
    \begin{itemize}
        \item Most RL algorithms aim to estimate the State(-Action) Value, i.e., how valuable is it to be in a certain state
        \item Value-based algorithms directly use this estimate to derive a policy
        \item Q-learning \liturl{Watkins, 1989}{https://www.cs.rhul.ac.uk/~chrisw/new_thesis.pdf} is one of the most popular value-based RL algorithms:
    \end{itemize}
    \begin{equation*}
        \mathcal{Q}(s_t, a_t) \gets \mathcal{Q}(s_t, a_t) +
        \alpha
        \underbrace{
            \Big(
                \overbrace{
                    r_{t} + \gamma \max\mathcal{Q}(s_{t+1},\cdot)
                }^{\text{TD-target}}
                - \mathcal{Q}(s_t,a_t)
            \Big),
        }_{\text{TD-delta}}
    \end{equation*}
\end{frame}
%-----------------------------------------------------------------------
%----------------------------------------------------------------------
\begin{frame}{Q-Learning Cont'd}
    \begin{equation*}
        \mathcal{Q}(s_t, a_t) \gets \mathcal{Q}(s_t, a_t) +
        \alpha
        \underbrace{
            \Big(
                \overbrace{
                    r_{t} + \gamma \max\mathcal{Q}(s_{t+1},\cdot)
                }^{\text{TD-target}}
                - \mathcal{Q}(s_t,a_t)
            \Big),
        }_{\text{TD-delta}}
    \end{equation*}
    \begin{description}
        \item[TD-target] Discounted expected future rewards given the assumption that we will behave greedy in the future. The discounting factor $\gamma$ determines how strongly to weigh the influence of predicted future rewards.
        In supervised learning terms: Defines the regression target value\pause
        \item[TD-delta] Difference between predicted reward $\mathcal{Q}(s_t,a_t)$ and the actual observed $r_t$ with expected future rewards.
        In supervised learning terms. Defines the prediction error
    \end{description}\pause
    The learning rate $\alpha$ determines how strongly to overwrite the current estimate.
\end{frame}
%-----------------------------------------------------------------------
%----------------------------------------------------------------------
\begin{frame}{The Policy of a Q-Learning Agent}
    \centering
    \begin{equation*}
        \mathcal{Q}(s_t, a_t) \gets \mathcal{Q}(s_t, a_t) +
        \alpha
        \underbrace{
            \Big(
                \overbrace{
                    r_{t} + \gamma \max\mathcal{Q}(s_{t+1},\cdot)
                }^{\text{TD-target}}
                - \mathcal{Q}(s_t,a_t)
            \Big),
        }_{\text{TD-delta}}
    \end{equation*}
    The policy then simply is $\policy(s)=\argmax_{a}\mathcal{Q}(s,\cdot)$
\end{frame}
%-----------------------------------------------------------------------
%----------------------------------------------------------------------
% \begin{frame}{Main Takeaways} 

% %Most important insights from this video:

% \begin{itemize}
%     \item
% \end{itemize}

% \end{frame}


	

% %----------------------------------------------------------------------
% %----------------------------------------------------------------------
% \begin{frame}[c]{Planning Slide}
% \begin{enumerate}
%     \item General Intro (MDPs, sequential decision making, policy learning)
%     \item Maybe Intro to AutoRL??? (only iff time permits)
% \end{enumerate}
% \end{frame}
%-----------------------------------------------------------------------

\end{document}
